\documentclass[aspectratio=169,10pt]{beamer}
\usetheme[numbering=fraction,progressbar=frametitle]{metropolis}
\usepackage{fontspec}
\setsansfont{Latin Modern Sans}
\usepackage{amsmath,amssymb,mathtools,graphicx,booktabs}
\usepackage{polyglossia}
\setdefaultlanguage{french}
\setbeamertemplate{navigation symbols}{}
\setbeamerfont{frametitle}{size=\large}
\setbeamerfont{normal text}{size=\normalsize}
\setlength{\parskip}{.45em}
\hypersetup{pdfauthor={},pdfsubject={Lecture slides},pdfkeywords={},pdftitle={Logique et raisonnement}}
\providecommand{\N}{\mathbb{N}}
\providecommand{\Z}{\mathbb{Z}}
\providecommand{\D}{\mathbb{D}}
\providecommand{\Q}{\mathbb{Q}}
\providecommand{\R}{\mathbb{R}}
\providecommand{\C}{\mathbb{C}}
\providecommand{\vect}[1]{\overrightarrow{#1}}
\providecommand{\vi}{\vec{\imath}}
\providecommand{\vj}{\vec{\jmath}}
\providecommand{\vk}{\vec{k}}
\providecommand{\norm}[1]{\left\lVert #1\right\rVert}
\providecommand{\vc}[2]{\begin{pmatrix}#1\\#2\end{pmatrix}}
\providecommand{\deter}[2]{\det\left(#1,#2\right)}
\providecommand{\abs}[1]{\left\lvert #1\right\rvert}
\providecommand{\dd}{\,\mathrm{d}}
\providecommand{\e}{\mathrm{e}}
\providecommand{\Card}{\operatorname{Card}}
\providecommand{\E}{\operatorname{E}}
\providecommand{\V}{\operatorname{V}}
\providecommand{\degre}{^{\circ}}
\providecommand{\un}[1]{\,\mathrm{#1}}
\providecommand{\sci}[2]{#1\times10^{#2}}
\providecommand{\pgcd}{\operatorname{PGCD}}
\providecommand{\ppcm}{\operatorname{PPCM}}
\title{Logique et raisonnement}
\subtitle{Cours et exemples}
\date{}
\author{}
\begin{document}
\begin{frame}[plain]
\vfill
{\usebeamerfont{title}\inserttitle\par}
\medskip
{\usebeamerfont{subtitle}\insertsubtitle\par}
\medskip
{\usebeamercolor[fg]{progress bar}\rule{\textwidth}{.5pt}}
\vfill
\end{frame}
\begin{frame}{Objectifs}
Formuler et nier des propositions, distinguer réciproque et contraposée et construire une démonstration ou un contre-exemple.
\end{frame}
\begin{frame}{Propositions et connecteurs}
Une proposition est un énoncé ayant une valeur de vérité. \(P\wedge Q\) signifie « \(P\) et \(Q\) »; \(P\vee Q\) signifie « \(P\) ou \(Q\) », au moins une étant vraie (ou inclusif). \(\neg P\) est la négation. L’implication \(P\Rightarrow Q\) est fausse seulement lorsque \(P\) est vraie et \(Q\) fausse. L’équivalence \(P\Leftrightarrow Q\) signifie les deux implications.
\end{frame}
\begin{frame}{Propositions et connecteurs}
\(\neg(P\wedge Q)\Leftrightarrow(\neg P\vee\neg Q)\);
\(\neg(P\vee Q)\Leftrightarrow(\neg P\wedge\neg Q)\).
La contraposée de \(P\Rightarrow Q\) est \(\neg Q\Rightarrow\neg P\) : elle est équivalente à l’implication initiale. La réciproque \(Q\Rightarrow P\) ne l’est pas en général.
\end{frame}
\begin{frame}{Propositions et connecteurs}
Les deux implications \(P\Rightarrow Q\) et \(\neg Q\Rightarrow\neg P\) sont fausses exactement dans le cas \(P\) vraie, \(Q\) fausse; elles ont donc la même table. Les lois de négation se vérifient sur les quatre lignes.
\end{frame}
\begin{frame}{Propositions et connecteurs}
« Si un entier est divisible par \(6\), il est divisible par \(3\) » est vrai. Sa réciproque est fausse : \(3\) est divisible par \(3\), mais pas par \(6\). Sa contraposée, « s’il n’est pas divisible par \(3\), il n’est pas divisible par \(6\) », est vraie.
\end{frame}
\begin{frame}{Quantifier et nier}
\(\forall x\in E\) signifie « pour tout \(x\) de \(E\) »; \(\exists x\in E\) signifie « il existe au moins un \(x\) de \(E\) ». Le domaine \(E\) fait partie de l’énoncé.
\end{frame}
\begin{frame}{Quantifier et nier}
\(\neg(\forall x\in E, P(x))\Leftrightarrow\exists x\in E,\neg P(x)\).
De même, \(\neg(\exists x\in E,P(x))\Leftrightarrow\forall x\in E,\neg P(x)\).
\end{frame}
\begin{frame}{Quantifier et nier}
Dire que tous les éléments vérifient \(P\) est faux exactement lorsqu’au moins un élément ne la vérifie pas. Dire qu’il en existe un est faux lorsque tous échouent.
\end{frame}
\begin{frame}{Quantifier et nier}
\(\forall x\in\mathbb{R},\exists y\in\mathbb{R},y>x\) est vrai : choisir \(y=x+1\). \(\exists y\in\mathbb{R},\forall x\in\mathbb{R},y>x\) est faux : pour le \(y\) proposé, prendre \(x=y\). L’ordre des quantificateurs change le sens.
\end{frame}
\begin{frame}{Choisir une preuve}
Preuve directe : partir des hypothèses et atteindre la conclusion. Contraposée : supposer la conclusion fausse et montrer l’hypothèse fausse. Absurde : supposer l’hypothèse vraie et la conclusion fausse, puis obtenir une contradiction. Pour réfuter une affirmation universelle, un contre-exemple suffit; pour la prouver, des exemples ne suffisent pas.
\end{frame}
\begin{frame}{Choisir une preuve}
Si \(n=2k\) avec \(k\in\mathbb{Z}\), alors \(n^2=4k^2=2(2k^2)\) est pair. Pour la réciproque, si \(n\) est impair, \(n=2k+1\), et \(n^2=2(2k^2+2k)+1\) est impair : contraposée.
\end{frame}
\begin{frame}{Choisir une preuve}
Dans \(P\Rightarrow Q\), vérifier \(P\) suffit pour conclure \(Q\); vérifier \(Q\) ne permet pas en général de remonter à \(P\). On dit que \(P\) est une condition suffisante de \(Q\), et \(Q\) une condition nécessaire de \(P\). Pour une équivalence, prouver séparément les deux sens, sauf si chaque transformation intermédiaire est elle-même équivalente.
\end{frame}
\begin{frame}{Choisir une preuve}
Être divisible par \(4\) suffit pour qu’un entier soit pair. Être pair est nécessaire pour être divisible par \(4\), mais ne suffit pas : \(6\) est pair sans être divisible par \(4\). Un seul contre-exemple réfute la réciproque; dix exemples divisibles par \(4\) ne la prouvent pas.
\end{frame}
\begin{frame}{Choisir une preuve}
Fixer le domaine; remplacer chaque « pour tout » par « il existe » et inversement, dans l’ordre; nier le prédicat final. La négation de \(<\) est \(\ge\), celle de « et » est « ou ». Une négation affirme exactement que l’énoncé initial est faux, sans modifier ses hypothèses ou son domaine.
\end{frame}
\begin{frame}{Choisir une preuve}
La négation de \(\forall x\in\mathbb{R},\exists y\in\mathbb{R},y^2=x\) est \(\exists x\in\mathbb{R},\forall y\in\mathbb{R},y^2\ne x\). La première est fausse et la seconde vraie : prendre \(x=-1\), car tous les carrés réels sont positifs ou nuls. Écrire seulement « il existe \(y\) tel que \(y^2\ne x\) » aurait laissé un autre sens.
\end{frame}
\end{document}
